\chapter{Stage 1: Pre-Routing Confidence ($u_{\text{pre}}$)}
\label{ch:upre}

We now open the first box of the per-query pipeline. Before \caem{} retrieves
anything, generates anything, or verifies anything, it asks a single cheap
question about the query in front of it: how ready is the model to answer this on
its own? The answer is one number, $\upre$, the pre-routing confidence. It is the
earliest signal in the pipeline and the only one computed before the system
commits to a tier, which is exactly why it carries a special responsibility,
namely the safety override that refuses to let an unready model answer without
grounded retrieval.

\begin{intuition}[A glance before the work]
Imagine a triage nurse who, in a few seconds and without running any tests, has to
decide whether a patient can wait or needs immediate care. The nurse is not
diagnosing; the nurse is deciding \emph{how much machinery to spend}. Stage 1 is
that glance. From one short look at the query, before any real answer exists, it
produces a confidence that decides whether the model has earned the right to take
a cheap path or must be sent down the expensive, grounded one. It is deliberately
fast and deliberately humble: when in any doubt, it escalates.
\end{intuition}

\section{What Stage 1 is for}
\label{sec:upre-job}

It helps to be exact about the one job $\upre$ does, because it is narrower than
its prominence suggests. The pre-routing confidence has a single decisive
consumer in the router of \Cref{ch:router}: the \textbf{safety override}. The
router checks, before anything else, whether $\upre$ has cleared a fixed floor. If
it has not, the query is forced off the cheap paths and onto the
retrieval-augmented tier, no matter how good a memory match might look.

\begin{pitfall}[$\upre$ does not enter the dispatch score]
It is natural to assume that a confidence computed this early feeds into the
tier-selection score. In the live code it does not. The router's combined score
is built only from the memory similarity and the neighbour's stored quality, the
subject of \Cref{ch:router}. The pre-routing confidence enters the router at
exactly one point, the safety override that is evaluated \emph{before} the score,
and nowhere else. Keeping these two mechanisms separate is deliberate: the score
answers ``does memory cover this query?'', while $\upre$ answers a different
question, ``is the model ready to generate at all?'' Conflating them would
over-retrieve when memory is strong but the model is shaky, and under-retrieve in
the opposite case.
\end{pitfall}

So $\upre$ is the model's pre-generation readiness, distilled into one calibrated
scalar, and its power is the veto. The rest of this chapter builds that scalar from
two cheap signals, fuses them, calibrates the result, and then shows the veto in
action.

\section{Signal one: token confidence ($u_{\text{token}}$)}
\label{sec:upre-utoken}

The first ingredient is the model's own commitment to its words. We met the raw
material in \Cref{sec:prereq-llm}: at each step the model assigns a probability to
the token it actually emits. If those probabilities are high, the model is moving
through the answer decisively; if they are low, it is hedging across many
alternatives at every step, which is the textual signature of guessing.

To turn this into a signal, Stage 1 runs a \emph{short} greedy decode, just
thirty-two tokens, under the chat template, and records the probability the model
assigned to each chosen token. It then takes the geometric mean of those
probabilities, which is the same as exponentiating the average log-probability.
The geometric mean is the right average here because probabilities multiply along
a sequence; averaging in log-space keeps one very low token from being masked by
many confident ones.

\begin{precise}[The $u_{\text{token}}$ definition]
Let the short decode produce tokens $\hat v_1, \dots, \hat v_n$ (stopping at the
end-of-sequence token), each with model probability $p_t(\hat v_t)$. Then
\[
  \utoken
  \;=\;
  \exp\!\left(\frac{1}{n}\sum_{t=1}^{n}\log p_t(\hat v_t)\right)
  \;=\;
  \Big(\textstyle\prod_{t=1}^{n} p_t(\hat v_t)\Big)^{1/n}.
\]
This is the per-token average probability under the model: a value near one means
the model was confident at every step, a value near zero means it was unsure
somewhere along the way. The decode is capped at thirty-two tokens because Stage 1
needs a confidence proxy, not a full answer, and the short cap is what keeps the
whole stage inside its latency budget.
\end{precise}

The actual computation is short enough to read in full, and it is worth doing so,
because two details in it matter.

\begin{lstlisting}[caption={The $u_{\text{token}}$ computation, condensed from \texttt{caem/confidence/pre\_routing.py}.}, label={lst:utoken}]
out = model.generate(**inputs, max_new_tokens=32, do_sample=False,
                     output_scores=True, return_dict_in_generate=True)

log_probs = []
for t, step_logits in enumerate(out.scores):          # one entry per generated token
    log_soft  = log_softmax(step_logits[0], dim=-1)
    token_id  = out.sequences[0, gen_start + t].item()
    if token_id == tokenizer.eos_token_id:
        break                                          # stop at end-of-sequence
    log_probs.append(log_soft[token_id].item())        # log P(chosen token)

u_token = exp(sum(log_probs) / len(log_probs))         # geometric mean
\end{lstlisting}

The first detail is the greedy, deterministic decode (\code{do\_sample=False}):
the signal must be reproducible, so the same query always yields the same
$\utoken$. The second is the fail-safe. If generation throws, or returns no
scorable tokens, the function returns $0.0$ rather than raising. A zero
pre-routing confidence is the most cautious value possible, and as we will see it
forces the query onto the grounded tier. Failure in Stage 1 therefore degrades
toward safety, never toward a confident wrong answer.

\begin{pitfall}[Raw text collapses the signal; the chat template rescues it]
One non-obvious requirement: the query must be wrapped in the chat-markup envelope
before the decode. An instruction-tuned model fed a bare question scatters its
probability mass across the vocabulary, and $\utoken$ collapses to roughly
$10^{-9}$ for every query, informative about nothing. Wrapping the query as a user
turn with a generation prompt puts the model into ``assistant answering a user''
mode, where its per-token probabilities become meaningful. The envelope adds about
twenty fixed tokens, identical for every query, which is why it does not distort
the second signal we turn to next.
\end{pitfall}

\section{Signal two: internal convergence ($c_{\text{conv}}$)}
\label{sec:upre-cconv}

Token confidence reads the model's output. The second signal reads its
\emph{internals}, and asks a different question: has the model settled on a single,
stable interpretation of the query, or is it still of several minds about what is
being asked? This is the internal-convergence ratio, adapted
from~\cite{nandakishor2025routing}.

The idea rests on how a decoder processes a query. The same query produces a stack
of hidden-state representations, one per layer. When the model has a clear,
unambiguous reading of the question, the representation stabilises: the later
layers vary little, having converged on an interpretation. When the question is
ambiguous or unfamiliar, the representation keeps shifting across layers, and the
early-layer variance stays high relative to the late layers. The ratio of those
two variances is a cheap proxy for ``did the model converge?''

\begin{figure}[t]
\centering
\begin{tikzpicture}[font=\small\sffamily]
  % left: not converged (high early variance)
  \begin{scope}[xshift=0cm]
    \node[font=\bfseries\sffamily, text=cbPitfall!75!black] at (1.6,4.3) {Unstable reading};
    \node[font=\scriptsize, text=cbInk!70] at (1.6,3.95) {early layers scattered $\Rightarrow$ high $c_{\text{conv}}$};
    % early layers (scattered dots)
    \foreach \y/\jit in {3.2/0.5, 2.9/-0.4, 3.05/0.35, 2.75/-0.55, 3.25/0.15}
      \filldraw[cbPitfall!70] (1.6+\jit,\y) circle (1.6pt);
    \node[font=\scriptsize,text=cbInk!60,anchor=west] at (2.5,3.0) {early};
    \draw[cbInk!25, dashed] (0.2,2.5) -- (3.0,2.5);
    % late layers (still somewhat scattered)
    \foreach \y/\jit in {1.9/0.25, 1.7/-0.2, 1.8/0.1, 1.6/-0.3, 2.0/0.18}
      \filldraw[cbPitfall!40] (1.6+\jit,\y) circle (1.6pt);
    \node[font=\scriptsize,text=cbInk!60,anchor=west] at (2.5,1.8) {late};
    \draw[->,cbInk!50] (0.2,1.2) -- (0.2,3.6) node[above,font=\scriptsize]{layer};
  \end{scope}
  % right: converged (low late variance)
  \begin{scope}[xshift=6.2cm]
    \node[font=\bfseries\sffamily, text=cbExample!60!black] at (1.6,4.3) {Stable reading};
    \node[font=\scriptsize, text=cbInk!70] at (1.6,3.95) {late layers tight $\Rightarrow$ low $c_{\text{conv}}$};
    \foreach \y/\jit in {3.2/0.45, 2.9/-0.35, 3.05/0.3, 2.75/-0.5, 3.25/0.1}
      \filldraw[cbExample!70] (1.6+\jit,\y) circle (1.6pt);
    \node[font=\scriptsize,text=cbInk!60,anchor=west] at (2.5,3.0) {early};
    \draw[cbInk!25, dashed] (0.2,2.5) -- (3.0,2.5);
    \foreach \y/\jit in {1.78/0.05, 1.74/-0.04, 1.8/0.02, 1.72/-0.05, 1.79/0.03}
      \filldraw[cbExample!70] (1.6+\jit,\y) circle (1.6pt);
    \node[font=\scriptsize,text=cbInk!60,anchor=west] at (2.5,1.78) {late};
    \draw[->,cbInk!50] (0.2,1.2) -- (0.2,3.6) node[above,font=\scriptsize]{layer};
  \end{scope}
\end{tikzpicture}
\caption[Internal convergence]{The convergence ratio compares the spread of the
early-layer representations with that of the late layers. On the left the late
layers are still scattered: the model has not converged on one reading, the ratio
$c_{\text{conv}}$ is high, and the confidence $1/(1+c_{\text{conv}})$ is low. On
the right the late layers are tight: the model has settled, $c_{\text{conv}}$ is
low, and the confidence is high. (Schematic.)}
\label{fig:cconv}
\end{figure}

\begin{precise}[The $c_{\text{conv}}$ definition]
A single forward pass over the query, with no generation, exposes the hidden-state
stack $\mathbf{H}^{(1)}, \dots, \mathbf{H}^{(L)}$ for the $L$ decoder layers (the
backbone has $L=36$). Splitting the stack into an early half and a late half, the
convergence ratio is the variance of the early half divided by that of the late
half,
\[
  \cconv
  \;=\;
  \frac{\operatorname{Var}\!\big(\mathbf{H}^{(1:L/2)}\big)}
       {\operatorname{Var}\!\big(\mathbf{H}^{(L/2+1:L)}\big) + \varepsilon},
\]
and it is turned into a confidence by the decreasing map
\[
  \text{confidence from } \cconv \;=\; \frac{1}{1 + \cconv}.
\]
A high ratio (early layers far more variable than late) reads as instability and
maps toward zero; a low ratio reads as a settled representation and maps toward
one. As with $\utoken$, any failure in this forward pass returns a ratio of zero,
the most confident value, but here the fail-safe is overridden in fusion because
the token signal will usually have failed too, and the combined floor still
escalates. The embedding layer is excluded from both halves so that only the
transformer computation, not the raw lookup, is measured.
\end{precise}

\section{Fusing and calibrating the two signals}
\label{sec:upre-fusion}

Two signals become one through a fixed weighted blend, followed by a calibration
step. The blend places more weight on the token signal, which is the more direct
reliability cue, and uses the convergence signal as a secondary prior on query
stability.

\begin{precise}[From two signals to a calibrated $u_{\text{pre}}$]
The raw pre-routing confidence is the registered blend
\[
  \upre^{\text{raw}}
  \;=\;
  0.60 \cdot \utoken \;+\; 0.40 \cdot \frac{1}{1 + \cconv},
\]
clipped to $[0,1]$. The mixing weights $0.60$ and $0.40$ are \emph{registered
before} the calibration phase rather than learned, so the meaning of the fused
scalar does not drift as the generator is fine-tuned across cycles. The raw value
is then passed through temperature scaling (\Cref{sec:prereq-calibration}) to make
it a calibrated probability,
\[
  \upre
  \;=\;
  \sigma\!\left(\frac{\operatorname{logit}\big(\upre^{\text{raw}}\big)}{T_b}\right),
\]
where the temperature $T_b$ is fit \emph{per benchmark}. A temperature of one
leaves the value untouched; a larger temperature softens an overconfident raw
distribution toward the diagonal of the reliability diagram.
\end{precise}

\begin{figure}[t]
\centering
\begin{tikzpicture}[
  font=\small\sffamily,
  sig/.style={rectangle, rounded corners=4pt, draw=cbIntuition!75!black,
    fill=cbIntuition!10, line width=1pt, align=center, inner sep=5pt,
    text width=3.5cm, minimum height=1.1cm, drop shadow={shadow xshift=0.3ex,
    shadow yshift=-0.3ex, fill=black!10}},
  op/.style={rectangle, rounded corners=4pt, draw=cbFormula!75!black,
    fill=cbFormula!10, line width=1pt, align=center, inner sep=5pt,
    text width=2.5cm, minimum height=1.1cm, drop shadow={shadow xshift=0.3ex,
    shadow yshift=-0.3ex, fill=black!10}},
  dec/.style={diamond, aspect=2.2, draw=cbPitfall!75!black, fill=cbPitfall!10,
    line width=1pt, align=center, inner sep=1pt, font=\footnotesize,
    text width=2.4cm},
  cbout/.style={rectangle, rounded corners=2pt, draw=cbInk!55, fill=cbInk!6,
    line width=1pt, align=center, inner sep=5pt},
  flow/.style={-{Stealth[length=2.4mm]}, line width=1.1pt, draw=cbInk!75},
  veto/.style={-{Stealth[length=2.4mm]}, line width=1.1pt, draw=cbPitfall!75, dashed},
]
  \node[sig] (ut) at (0,1.1) {\textbf{32-token greedy decode}\\$\utoken$ (geom.\ mean)};
  \node[sig] (cc) at (0,-1.1) {\textbf{one forward pass}\\$1/(1+\cconv)$};
  \node[op]  (blend) at (4.0,0) {\textbf{blend}\\$0.6\,\utoken + 0.4(\cdot)$};
  \node[op]  (temp) at (7.3,0) {\textbf{temperature}\\$\sigma(\operatorname{logit}/T_b)$};
  \node[dec] (gate) at (10.7,0) {$\upre \ge 0.38$?};
  \node[cbout] (proceed) at (10.7,1.9) {to dispatch score (\Cref{ch:router})};
  \node[cbout, text=cbPitfall!70!black] (t3) at (10.7,-1.9) {\textbf{safety override $\to$ Tier 3}};

  \draw[flow] (ut) -- (blend);
  \draw[flow] (cc) -- (blend);
  \draw[flow] (blend) -- (temp);
  \draw[flow] (temp) -- (gate);
  \draw[flow] (gate) -- node[right,font=\scriptsize,text=cbExample!55!black]{yes} (proceed);
  \draw[veto] (gate) -- node[right,font=\scriptsize,text=cbPitfall!70!black]{no} (t3);
\end{tikzpicture}
\caption[The Stage 1 data flow]{The pre-routing confidence is built from two cheap
signals, a thirty-two-token greedy decode for $\utoken$ and a single forward pass
for $\cconv$, blended with fixed weights, calibrated by a per-benchmark
temperature, and then tested against the safety floor. Clearing the floor hands
control to the dispatch score; failing it triggers the safety override straight to
the grounded tier. (Schematic.)}
\label{fig:upre-flow}
\end{figure}

\section{The safety override}
\label{sec:upre-override}

Now the payoff. The router evaluates one rule before any other: if $\upre$ is
below the registered floor $\phisafe = 0.38$, the query is dispatched to the
retrieval-augmented tier and the override flag is recorded, regardless of how
strong the memory match is. This is the architectural expression of a single
principle: \textbf{the system must not generate without grounded retrieval
whenever its pre-generation readiness is low}. The cost of grounded retrieval is
explicitly preferred over the risk of a confident wrong answer from memory or
zero-shot generation.

\begin{fromcode}[The floor, and why it is $0.38$ rather than $0.60$]
The floor began life at $0.60$ and was lowered to $0.38$ during the cycle-one
routing audit. The audit found that the higher floor was forcing well-supported
queries onto the grounded tier under the post-calibration pre-routing
distribution, which depressed the architecture's effective use of the cheaper
tiers without a matching gain in safety. Lowering the floor preserved the
override's intent, refuse to generate when readiness is genuinely low, while
admitting a wider band of capable queries to the cheap paths. The floor is read
per benchmark through a configuration helper, with the pooled $\phisafe = 0.38$ used
as the registered default. The per-benchmark floors are re-fit across cycles, and
they are not all equal: the fact-checking and trivia probes hold at the pooled $0.38$,
while the commonsense probe drifts upward, to about $0.49$ by the third cycle and
about $0.57$ by the fifth, as its calibrated readiness distribution shifts. The pooled
$0.38$ remains the fallback wherever a benchmark-specific value is unavailable.
\end{fromcode}

The ordering is the subtle part. Because the override is checked first, a query
with a near-perfect memory match but a low readiness signal is still sent to
retrieval. The override is a veto, not a vote: it cannot be outweighed by a good
score, only avoided by clearing the floor. When the floor is cleared, $\upre$ has
done its job and steps aside; the tier is then chosen by the similarity-and-quality
score of \Cref{ch:router}.

On the realised trajectory the override is a worst-case safety contract that, in
fact, never activates: the calibrated pre-routing confidence stays above the active
floor on every one of the roughly nine thousand evaluation queries, so the override
never fires. This is best read not as evidence that the mechanism is
unneeded, but as evidence that the panel distribution simply does not exercise it. The
contract is cheap to carry and would fire the moment a genuinely unready query
appeared; the deployed query mix never produced one that fell through the floor.

\section{Why Stage 1 is cheap by design}
\label{sec:upre-cost}

A theme of the whole architecture is paying for confidence only where it is
needed, and Stage 1 is the first instance. The pre-routing confidence pointedly
does \emph{not} use the expensive multi-chain resampling that the post-generation
verifier of \Cref{ch:verifier} relies on. It costs one forward pass over the query
plus a thirty-two-token greedy decode, and that is all.

\begin{intuition}[The right signal at the right price]
There are two confidence questions in the pipeline, asked at different prices.
Stage 1 asks the cheap one, ``should the model even try?'', and pays a single
short decode for it, because every query must be triaged. The expensive question,
``is this generated answer trustworthy enough to keep?'', is asked only of queries
that actually reach the generating tiers, and is paid for with the resampling and
entailment machinery of \Cref{ch:verifier}. A memory-direct query never pays the
expensive cost at all. Spending a cheap signal on everyone and an expensive signal
on the few who need it is the economic logic of the entire compute ladder.
\end{intuition}

\section{The two readings, at the $u_{\text{pre}}$ level}
\label{sec:upre-scenarios}

The two queries we traced in \Cref{ch:bigpicture} diverge for the first time
exactly here, at Stage 1, and it is instructive to watch only this stage.

\begin{scenario}[The claim that clears the floor (FEVER)]
For the claim ``The French Revolution began before 1792'', the short decode commits
to its tokens with reasonable probability and the forward pass shows a settled
late-layer representation. The blended raw confidence, after the per-benchmark
temperature, lands comfortably above $0.38$. The safety override does \emph{not}
fire. Control passes to the dispatch score, and because the cycle-zero memory is
empty, the query continues down the pipeline toward generation rather than being
vetoed here. Stage 1's verdict: this query is ready enough; let the later stages
decide the tier.
\end{scenario}

\begin{scenario}[The question routed to grounding (trivia)]
For ``\,`If I Ruled The World' comes from which musical?'', the short decode commits
only weakly to its tokens (the token-probability aggregate sits near $0.02$), yet the
calibrated pre-routing confidence still lands around $\upre = 0.66$, comfortably above
the $0.38$ safety floor. The override therefore does \emph{not} fire here. The query
nonetheless goes to the retrieval-augmented tier, because at the cold start the
episodic memory is empty and no cheap path can be earned: the dispatch score finds no
neighbour to draw on. Grounded generation is what later gives the verifier something
real to judge, and ultimately what lets the system refuse rather than fabricate. The
honest refusal at the end of the pipeline begins with this early decision to ground.
\end{scenario}

\section{An honest limitation}
\label{sec:upre-threats}

One property of $\upre$ deserves a frank note, because it surfaces in the
trajectory of \Cref{ch:trajectory}. The single-parameter temperature that calibrates
the pre-routing confidence (distinct from the composite's own temperature in
\Cref{ch:composite}) behaves oddly across the run. At the cold start the fit saturates
its optimisation bound, returning a temperature near $20$ that is best flagged as
pathological calibration data; even there the calibration error does fall, from about
$0.28$ to about $0.07$. The per-cycle re-fits that follow then settle at low interior
values, roughly $1.09$, $0.84$, and $0.74$ over the next cycles, while the calibration
error barely moves, staying around $0.04$ to $0.05$. The honest reading is therefore
not that the temperature is pinned at its ceiling every cycle, but that
single-parameter temperature scaling on this particular surface is largely inert
across the trajectory: it neither helps nor hurts $\upre$ much once the cold-start
boundary is past.

\begin{pitfall}[Why the drift is tolerable]
A miscalibrated $\upre$ would be a serious problem if it gated what \caem{}
stores. It does not. The pre-routing confidence governs only \emph{routing}: the
safety override and, indirectly, which tier pays the verification cost. The
decision of whether an answer is trustworthy enough to keep is made downstream, by
the calibrated composite of \Cref{ch:composite}, on signals that carry their own
per-cycle calibration. So a drift in $\upre$ can send a few queries down a more
expensive tier than strictly necessary, which costs compute, but it cannot admit a
wrong answer into memory. The architecture localises the consequences of this
imperfection to cost, not correctness, which is precisely why the cheap signal is
allowed to be imperfect.
\end{pitfall}

\begin{bythenumbers}[Stage 1 constants]
\begin{itemize}[leftmargin=1.3em, itemsep=1pt]
  \item $\utoken$ decode length: $32$ tokens, greedy, deterministic.
  \item Fusion weights: $0.60$ on $\utoken$, $0.40$ on $1/(1+\cconv)$, registered (not learned).
  \item Convergence split: early half versus late half of $L=36$ decoder layers, embedding layer excluded.
  \item Calibration: per-benchmark temperature $T_b$ via $\sigma(\operatorname{logit}(\cdot)/T_b)$.
  \item Safety floor: $\phisafe = 0.38$ (lowered from $0.60$ at the cycle-one audit).
  \item Fail-safe: any error returns $0.0$, the most cautious value, forcing the grounded tier.
\end{itemize}
\end{bythenumbers}

\begin{defence}[If an examiner asks: ``is two signals not too few for a confidence estimate?'']
The pre-routing confidence is deliberately thin because of what it is for. It is a
triage signal whose only hard consequence is a conservative veto, backed by a
fail-safe that escalates on any doubt, so the cost of an occasional false
low-confidence reading is bounded compute rather than a wrong answer. The
literature is in any case discouraging about single-signal confidence, which
reaches only chance-to-modest discrimination on its own; \caem{} answers that not
by piling more signals into the cheap pre-generation estimate, but by deferring the
heavy, many-signal judgement to the post-generation verifier, where it can be
afforded. Two cheap signals to triage, nine calibrated signals to decide: the
split is the point, not an oversight.
\end{defence}

With the query triaged, the pipeline turns to memory. \Cref{ch:memory} takes up
Stage 2, the episodic store: how a query is encoded, how the nearest past episode
is found, and what an episode actually carries.
